\clearpage\SourceLeaf{464}
\IndexRow{0}{Zur Parallelentheorie:\quad \textsc{Schumacher} an \textsc{Gauss} 1831 Mai 2.}{.\,.\,.\,.\,.\,.\,.\,.\,.\,.\,.}{Seite}{210}
\IndexRow{\widthof{Zur Parallelentheorie:\quad}}{\textsc{Gauss} an \textsc{Schumacher} 1831 Mai 17.}{.\,.\,.\,.\,.\,.\,.\,.\,.\,.\,.}{—}{212}
\IndexRow{\widthof{Zur Parallelentheorie:\quad}}{\textsc{Schumacher} an \textsc{Gauss} 1831 Mai 25.}{.\,.\,.\,.\,.\,.\,.\,.\,.\,.\,.}{—}{213}
\IndexRow{\widthof{Zur Parallelentheorie:\quad}}{\textsc{Schumacher} an \textsc{Gauss} 1831 Juni 29.}{.\,.\,.\,.\,.\,.\,.\,.\,.\,.\,.}{—}{214}
\IndexRow{\widthof{Zur Parallelentheorie:\quad}}{\textsc{Gauss} an \textsc{Schumacher} 1831 Juli 12.}{.\,.\,.\,.\,.\,.\,.\,.\,.\,.\,.}{—}{215}
\IndexRow{\widthof{Zur Parallelentheorie:\quad}}{\textsc{Schumacher} an \textsc{Gauss} 1831 Juli 19.}{.\,.\,.\,.\,.\,.\,.\,.\,.\,.\,.}{—}{218}
\IndexRow{0}{\textsc{Johann Bolyais} Appendix:\quad \textsc{Gauss} an \textsc{Gerling} 1832 Febr. 14.}{.\,.\,.\,.\,.\,.\,.\,.\,.}{—}{220}
\IndexRow{\widthof{\textsc{Johann Bolyais} Appendix:\quad}}{\textsc{Gauss} an \textsc{W. von Bolyai} 1832 März 6.}{.\,.\,.\,.\,.\,.\,.}{—}{220}
\IndexRow{0}{Zur Astralgeometrie (\textit{Nachlass} 1832)}{.\,.\,.\,.\,.\,.\,.\,.\,.\,.\,.\,.\,.\,.\,.\,.\,.\,.\,.}{—}{226}
\IndexRow{0}{\textsc{Lübsens} Parallelentheorie:\quad \textsc{Schumacher} an \textsc{Gauss} 1835 Dec. 31.}{.\,.\,.\,.\,.\,.\,.\,.\,.}{—}{230}
\IndexRow{\widthof{\textsc{Lübsens} Parallelentheorie:\quad}}{\textsc{Gauss} an \textsc{Schumacher} 1836 Jan. 2.}{.\,.\,.\,.\,.\,.\,.\,.\,.}{—}{230}
\IndexRow{0}{Volumenbestimmungen in der Nichteuklidischen Geometrie (\textit{Brief} 1841 und \textit{Nachlass})}{.}{—}{232}
\IndexRow{0}{\textsc{Bolyai} und \textsc{Lobatschefsky} (\textit{Briefe} 1843—1846)}{.\,.\,.\,.\,.\,.\,.\,.\,.\,.\,.\,.\,.\,.\,.}{—}{234}
\IndexRow{0}{Congruenz und Symmetrie (\textit{Briefe} 1813 und 1844—1846)}{.\,.\,.\,.\,.\,.\,.\,.\,.\,.\,.}{—}{240}
\IndexRow{0}{Theorem aus der Sphärologie (\textit{Nachlass})}{.\,.\,.\,.\,.\,.\,.\,.\,.\,.\,.\,.\,.\,.\,.\,.\,.\,.}{—}{250}
\IndexRow{0}{Die sphärische und die Nichteuklidische Geometrie (\textit{Nachlass})}{.\,.\,.\,.\,.\,.\,.\,.\,.\,.\,.}{—}{255}
\IndexRow{0}{Über die Summe der Aussenwinkel eines Polygons:\quad \textsc{Gauss} an \textsc{Gerling} 1846 Oct. 2.}{.\,.\,.}{—}{266}
\IndexRow{0}{Metaphysik der Geometrie (\textit{Aus} \textsc{Sartorius von Waltershausen}, \textit{Gauss zum Gedächtniss})}{}{—}{267}
\IndexHeading{GEOMETRIA SITUS.}
\IndexSubheading{Nachlass.}
\IndexRow{0}{Zur Geometria situs}{.\,.\,.\,.\,.\,.\,.\,.\,.\,.\,.\,.\,.\,.\,.\,.\,.\,.\,.\,.\,.\,.\,.\,.\,.\,.}{—}{271}
\IndexRow{0}{Zur Geometrie der Lage, für zwei Raumdimensionen}{.\,.\,.\,.\,.\,.\,.\,.\,.\,.\,.\,.\,.}{—}{282}
\IndexHeading{AUFGABEN UND LEHRSÄTZE DER ELEMENTAREN GEOMETRIE ANGEHÖRIG.}
\IndexSubheading{Nachlass und Briefwechsel.}
\IndexRow{0}{Zur sphärischen Trigonometrie}{.\,.\,.\,.\,.\,.\,.\,.\,.\,.\,.\,.\,.\,.\,.\,.\,.\,.\,.\,.\,.\,.}{—}{289}
\IndexRow{0}{Geometrischer Ort der Spitze des sphärischen Dreiecks auf gegebener Basis mit gegebenem Inhalt}{.\,.\,.\,.\,.\,.\,.\,.\,.\,.\,.\,.\,.\,.\,.\,.\,.\,.\,.\,.\,.\,.\,.\,.\,.\,.\,.\,.\,.}{—}{292}
\IndexRow{0}{Zu \textsc{Möbius}' barycentrischem Calcul}{.\,.\,.\,.\,.\,.\,.\,.\,.\,.\,.\,.\,.\,.\,.\,.\,.}{—}{295}
\IndexHeading{VERWENDUNG COMPLEXER GRÖSSEN FÜR DIE GEOMETRIE.}
\IndexSubheading{Nachlass und Briefwechsel.}
\IndexRow{0}{Das Dreieck}{.\,.\,.\,.\,.\,.\,.\,.\,.\,.\,.\,.\,.\,.\,.\,.\,.\,.\,.\,.\,.\,.\,.\,.\,.\,.\,.\,.}{Seite}{303}
\IndexRow{0}{\textsc{Pothenots} Aufgabe und das Viereck}{.\,.\,.\,.\,.\,.\,.\,.\,.\,.\,.\,.\,.\,.\,.\,.\,.\,.\,.}{—}{307}
\IndexRow{0}{Der Kreis}{.\,.\,.\,.\,.\,.\,.\,.\,.\,.\,.\,.\,.\,.\,.\,.\,.\,.\,.\,.\,.\,.\,.\,.\,.\,.\,.\,.\,.}{—}{335}
\IndexRow{0}{Die Kegelschnitte}{.\,.\,.\,.\,.\,.\,.\,.\,.\,.\,.\,.\,.\,.\,.\,.\,.\,.\,.\,.\,.\,.\,.\,.\,.\,.}{—}{339}
\IndexRow{0}{Projection des Würfels}{.\,.\,.\,.\,.\,.\,.\,.\,.\,.\,.\,.\,.\,.\,.\,.\,.\,.\,.\,.\,.\,.\,.\,.\,.}{—}{345}
\IndexRow{0}{Geometrische Seite der ternären Formen}{.\,.\,.\,.\,.\,.\,.\,.\,.\,.\,.\,.\,.\,.\,.\,.\,.\,.}{—}{348}
\IndexRow{0}{Die Kugel}{.\,.\,.\,.\,.\,.\,.\,.\,.\,.\,.\,.\,.\,.\,.\,.\,.\,.\,.\,.\,.\,.\,.\,.\,.\,.\,.\,.\,.}{—}{351}
\IndexRow{0}{Mutationen des Raumes}{.\,.\,.\,.\,.\,.\,.\,.\,.\,.\,.\,.\,.\,.\,.\,.\,.\,.\,.\,.\,.\,.\,.\,.}{—}{357}
